% =====================================================================
% B5 - COVERAGE TABLE
% Ready to \input. Requires \usepackage{booktabs} and \usepackage{amssymb}.
%
% Cite key assumed to be \cite{bbca} for Tas & Baktir, IEEE Access 2024.
% Adjust to match your .bib.
%
% NOTATION
%   \cmark  blocked / detected
%   \pmark  partially addressed, with a stated gap
%   \xmark  not addressed by this mechanism
%
% Add these to your preamble if not already present:
%   \newcommand{\cmark}{\ding{51}}
%   \newcommand{\xmark}{\ding{55}}
%   \newcommand{\pmark}{$\circ$}
%   \usepackage{pifont}
% =====================================================================

\begin{table*}[t]
\centering
\caption{Attack coverage by mechanism. Columns are the mechanisms in
isolation; the final column is BVI with all four layers composed under the
gating rule of Proposition~1. BBCA addresses caller-ID and ANI integrity only,
and permits writes solely from registered ISPs (Sec.~III.A,~\cite{bbca}), so a
path containing an unregistered hop yields no record at all. The digest and the
detector answer different questions and neither subsumes the other: the
detector asks whether speech was synthesised, the digest asks whether this is
the audio the counterparty sent. Rows~8 and~9 are the two attacks that motivate
composing the layers, since each defeats every mechanism except one.}
\label{tab:coverage}
\begin{tabular}{clcccc}
\toprule
& Attack & BBCA & Digest (L3) & Detector (L4) & BVI \\
\midrule

\multicolumn{6}{l}{\emph{Credential attacks --- Layer 1}} \\[2pt]

A1 & No credential presented
   & \xmark & \xmark & \xmark & \cmark \\
A2 & Self-signed credential
   & \xmark & \xmark & \xmark & \cmark \\
A3 & Stolen credential, valid signature
   & \pmark & \xmark & \xmark & \cmark \\
A4 & Tampered credential claim
   & \pmark & \xmark & \xmark & \cmark \\
A5 & Revoked credential
   & \xmark & \xmark & \xmark & \cmark \\
A6 & Expired credential / stale key epoch
   & \xmark & \xmark & \xmark & \cmark \\
A7 & Replayed session assertion
   & \xmark & \cmark & \xmark & \cmark \\

\addlinespace
\multicolumn{6}{l}{\emph{Path attack --- Layer 2}} \\[2pt]

A8 & In-path identifier rewrite, audio untouched
   & \pmark & \xmark & \xmark & \cmark \\

\addlinespace
\multicolumn{6}{l}{\emph{Content attacks --- Layers 3 and 4}} \\[2pt]

A9 & Media substitution in transit, route untouched
   & \xmark & \cmark & \xmark & \cmark \\
A10 & Synthesised (deepfake) speech
   & \xmark & \xmark & \cmark & \cmark \\

\bottomrule
\end{tabular}

\vspace{4pt}
\begin{minipage}{0.94\textwidth}
\footnotesize
\textbf{Notes on the partial marks.}
\emph{A3}: BBCA re-registers an ISP on every caller-ID change, which surfaces a
change but does not distinguish a legitimate rotation from a compromise, and
provides no revocation mechanism to terminate the credential once compromise is
known.
\emph{A4}: BBCA detects an ANI value that disagrees with its registration
table, but the table is the only referent; a claim consistent with a
compromised registration passes.
\emph{A8}: BBCA verifies at each hop against the Call Flow Control Table, so a
rewrite between two \emph{registered} hops is visible. A rewrite at or adjacent
to an unregistered hop is not, and the design defines no behaviour for that
case. BVI treats an unattested hop as $\bot$ rather than as failure, so
detection degrades gracefully rather than disappearing (Table~\ref{tab:participation}).
\end{minipage}
\end{table*}
